% Môn: Vật lý
% Lớp: 11
% Chương: Chương III. Điện trường
% Bài: L11C3  Bài 20. Điện thế
% Số câu: 9
% App đọc file này trực tiếp — sửa rồi Commit trên GitHub, bấm Đồng bộ GitHub .tex

% L11C3  Bài 20. Điện thế

% ===== Câu 1 =====
% ID: L11C3B20-01-TL
% Mức: NB
\dangbt{Chưa có dạng}
\begin{ex}
% ID: L11C3B20-01-TL
% Mức: NB
Để mô tả điện thế trong không gian người ta còn dùng các mặt đẳng thế. Mặt đẳng thế là các mặt được vẽ trong không gian sao cho điện thế của các điểm trên mặt đẳng thế là bằng nhau, vectơ pháp tuyến của mặt đẳng thế được chọn hướng theo chiều tăng của điện thế. a) Chứng minh rằng công của lực điện trong sự dịch chuyển các điện tích bên trong mặt đẳng thế luôn bằng 0. b) Chứng tỏ rằng vectơ pháp tuyến của mặt đẳng thế tại mỗi điểm cùng phương và ngược chiều với vectơ cường độ điện trường tại điểm đó. $\nguon{ly11 kntt.pdf}$
\choiceTF

\loigiai{
a) Với hai điểm lân cận $M$ và $N$ bất kì trên mặt đẳng thế, luôn có cùng điện thế nên công dịch chuyển một điện tích q giữa hai điểm luôn bằng không. Điều đó có nghĩa là vectơ lực điện $\overset\rightarrow{F}$ tác dụng lên một điện tích q nằm trên mặt đẳng thế sẽ vuông góc với mặt đẳng thế.

b) Chiều vectơ cường độ điện trường tại mỗi điểm trên mặt đẳng thế hướng theo chiều giảm của điện thế nên nó sẽ ngược chiều với vectơ pháp tuyến của mặt đẳng thế tại điểm đó.

Phương của vectơ cường độ điện trường tại mỗi điểm trên mặt đẳng thế cùng phương với vectơ lực điện $\overset\rightarrow{F}$ tác dụng lên một điện tích q tại điểm đó, tức là sẽ cùng phương với vectơ pháp tuyến của mặt đẳng thế tại điểm đó.

Kết luận: vectơ pháp tuyến của mặt đẳng thế tại mỗi điểm cùng phương và ngược chiều với vectơ cường độ điện trường tại điểm đó.
}
\end{ex}

% ===== Câu 2 =====
% ID: L11C3B20-01-TN
% Mức: NB
\begin{ex}
% ID: L11C3B20-01-TN
% Mức: NB
Đơn vị của điện thế là: $\nguon{ly11 kntt.pdf}$
\choice
{\True vôn (V).}
{jun (J).}
{vôn trên mét $(\mathrm{V/m})$.}
{oát (W).}
\loigiai{
Chọn A.

Đơn vị của điện thế là vôn (V).
}
\end{ex}

% ===== Câu 3 =====
% ID: L11C3B20-02-TL
% Mức: TH
\begin{ex}
% ID: L11C3B20-02-TL
% Mức: TH
Trong điện trường của một điện tích $Q$ cố định, công để dịch chuyển một điện tích q từ vô cùng $Q$ về điểm $M$ cách $Q$ một khoảng r có giá trị bằng $A$ \inftyM = q $.M$ là một điểm cách $Q$ một khoảng 4 $\pi\varepsilon_0$ r $1\,\text{m}$ và $N$ là một điểm cách $Q$ một khoảng $2\,\text{m}$. a) Hãy tính hiệu điện thế UM𝑁. b) Áp dụng với $Q = 8 \cdot 10^{-10} C$. Tính công cần thực hiện để dịch chuyển một electron từ $M$ đến N. $\nguon{ly11 kntt.pdf}$
\choiceTF

\loigiai{
a) $U_{MN} = \dfrac{A_{MN}}{q} = \dfrac{A_{M\infty} - A_{N\infty}}{q} = \dfrac{- q\dfrac{Q}{4\pi\varepsilon_{0}r_{M}} + q\dfrac{Q}{4\pi\varepsilon_{0}r_{N}}}{q} = - \dfrac{Q}{4\pi\varepsilon_{0}.1} + \dfrac{Q}{4\pi\varepsilon_{0}.2} = - \dfrac{Q}{8\pi\varepsilon_{0}}(V)$ b) $\begin{matrix}
{A_{MN} = A_{M\infty} - A_{N\infty} = - A_{\infty M} - \left( {- A_{\infty N}} \right) = - q\dfrac{Q}{4\pi\varepsilon_{0}r_{M}} - \left( {- q\dfrac{Q}{4\pi\varepsilon_{0}r_{N}}} \right)} \\
{= - q\dfrac{Q}{4\pi\varepsilon_{0}.1} + q\dfrac{Q}{4\pi\varepsilon_{0}.2} = - q\dfrac{Q}{8\pi\varepsilon_{0}} = \dfrac{1,6 \cdot 10^{-29}}{\pi\varepsilon_{0}}(J)}
\end{matrix}$
}
\end{ex}

% ===== Câu 4 =====
% ID: L11C3B20-02-TN
% Mức: NB
\begin{ex}
% ID: L11C3B20-02-TN
% Mức: NB
Điện thế tại một điểm $M$ trong điện trường bất kì có cường độ điện trường $E$ không phụ thuộc vào $\nguon{ly11 kntt.pdf}$
\choice
{vị trí điểm M.}
{cường độ điện trường E.}
{\True điện tích q đặt tại điểm M.}
{vị trí được chọn làm mốc của điện thế.}
\loigiai{
Chọn C.

Điện thế tại một điểm $M$ trong điện trường bất kì có cường độ điện trường $\overset\rightarrow{E}$ không phụ thuộc vào điện tích q đặt tại điểm M.
}
\end{ex}

% ===== Câu 5 =====
% ID: L11C3B20-03-TL
% Mức: TH
\begin{ex}
% ID: L11C3B20-03-TL
% Mức: TH
Một đám mây dông bị phân thành hai tầng, tầng trên mang điện dương cách xa tầng dưới mang điện âm. Đo bằng thực nghiệm, người ta thấy điện trường trong khoảng giữa hai tầng của đám mây dông đó gân đều, hướng từ trên xuống dưới với $E =$ 830 $\mathrm{V/m},$ khoảng cách giữa hai tầng là $0,7 \mathrm{km}$, điện tích của tầng phía trên ước tính được bằng $Q_1 = 1$, $24\,\mathrm{C}$. Coi điện thế của tầng mây phía dưới là V_1. a) Hãy tính điện thế của tầng mây phía trên. b) Ước tính thế năng điện của tầng mây phía trên. $\nguon{ly11 kntt.pdf}$
\choiceTF

\loigiai{
a) Vận dụng mối liên hệ giữa cường độ điện trường và điện thế $E_{M} = \dfrac{V_{M} - V_{N}}{\overline{MN}}$, với $M$ là điểm ở tầng mây phía dưới, $N$ là điểm ở tầng mây phía trên, $\overline{MN}$ được tính ngược chiều đường sức điện nên có giá trị âm, ta tính được điện thế của tầng mây phía trên: V_(N) = V₂ = V₁ - Eh = V₁ + 581000(V)

b) Vận dụng mối liên hệ giữa thế năng điện và điện thế W_(M) = Vq. Ta tính được thế năng điện của tầng mây phía trên: W₂ = (V₁+581000) ⋅ 1, 24(J)
}
\end{ex}

% ===== Câu 6 =====
% ID: L11C3B20-03-TN
% Mức: NB
\begin{ex}
% ID: L11C3B20-03-TN
% Mức: NB
Biết điện thế tại điểm $M$ trong điện trường đều trái đất là $120\,\mathrm{V}$. Mốc thế năng điện được chọn tại mặt đất. Electron đặt tại điểm $M$ có thế năng là: $\nguon{ly11 kntt.pdf}$
\choice
{$-192 \cdot 10^{-19} V.$}
{\True $-192 \cdot 10^{-19} J.$}
{$192 \cdot 10^{-19} V.$}
{$192 \cdot 10^{-19} J$}
\loigiai{
Chọn B.

Thế năng $A= qU = -1,6\cdot10^{-19}.120 = -1,92\cdot10^{-17}J$
}
\end{ex}

% ===== Câu 7 =====
% ID: L11C3B20-04-TL
% Mức: TH
\begin{bt}
% ID: L11C3B20-04-TL
% Mức: TH
Tiếp tục đo bằng thực nghiệm tầng mây phía dưới của đám mây dông ở Câu 7, người ta thấy nó nằm cách mặt đât khoảng $6450\,\text{m}$. Trong khoảng không gian nằm giữa mặt đất và tầng dưới đám mây có điện trường đều hướng thẳng đứng từ dưới lên trên với $E =$ 250 $\mathrm{V/m}$. Điện tích của tầng dưới đám mây ước tính được là $Q_2 = -2$, $03\,\mathrm{C}$. a) Chọn mốc điện thế là mặt đất, hãy ước tính điện thế của tầng phía dưới đám mây dông trên. b) Tính thế năng điện của tầng dưới đám mây dông đó. $\nguon{ly11 kntt.pdf}$
\loigiai{
a) Với chú ý điện trường hướng từ dưới lên trên Hình 20.1 $G$, $\overline{MN}$ được tính ngược chiều đường sức điện nên có giá trị âm, điểm $M$ là điểm ở tầng thấp đám mây, điểm $N$ là điểm trên mặt đất trong công thức $E_{M} = \dfrac{V_{M} - V_{N}}{\overline{MN}}$. Ta có điện thế của tầng thấp đám mây là: $V_{M} = V_{N} + E_{M}\overline{MN} = 0 + 250.\left( {- 6450} \right) = - 1612500\text{V}$

[Tiếp tục đo bằng thực nghiệm tầng mây phía dưới]

b) Thế năng điện của tầng dưới đám mây dông là:

W_(M) = Vq = -1612500.(-2,03) = $3273375\,\mathrm{J}$
}
\end{bt}

% ===== Câu 8 =====
% ID: L11C3B20-04-TN
% Mức: NB
\begin{ex}
% ID: L11C3B20-04-TN
% Mức: NB
Tại nơi có điện trường trái đất bằng $115 \mathrm{V/m}$, người ta đặt hai bản phẳng song song với nhau và song song với mặt đất. Bản thứ nhất cách mặt đất $1\,\text{m}$ và được nối với mặt đất bằng một dây đồng. Bản thứ hai cách mặt đất $1,073\,\text{m}$ và được tích điện dương. Hiệu điện thế đo được giữa hai bản là $1,5\,\mathrm{V}$. Chọn mặt đất là mốc điện thế, điện thế bản nhiễm điện dương bằng $\nguon{ly11 kntt.pdf}$
\choice
{\True $1,5\,\mathrm{V}$.}
{$8,39\,\mathrm{V}$.}
{$0\,\mathrm{V}$.}
{$-8,39\,\mathrm{V}$.}
\loigiai{
Chọn A.
}
\end{ex}

% ===== Câu 9 =====
% ID: L11C3B20-05-TN
% Mức: NB
\begin{ex}
% ID: L11C3B20-05-TN
% Mức: NB
Dọc theo đường sức điện của một điện tích âm được đặt trong chân không, điện thế sẽ $\nguon{ly11 kntt.pdf}$
\choice
{giảm dần khi đi từ điện tích ra xa vô cùng.}
{\True tăng dần khi đi từ điện tích ra xa vô cùng.}
{luôn không đổi vì các điểm nằm trên cùng một đường sức điện.}
{lúc đầu tăng lên sau đó giảm dần khi đi từ điện tích ra xa vô cùng.}
\loigiai{
Chọn B.

Dọc theo đường sức điện của một điện tích âm được đặt trong chân không, điện thế sẽ tăng dần khi đi từ điện tích ra xa vô cùng.
}
\end{ex}
